\section{Parallel execution and shared user-subroutine state}
\label{sec:externaldb-parallelization}

The use of \texttt{UEXTERNALDB} is not, by itself, incompatible with parallel
Abaqus/Standard execution.  The callback provides analysis- and
increment-boundary points at which external data can be initialized or
exchanged.  The relevant correctness question is instead how that information
is stored and subsequently accessed by element and material subroutines.

The present implementation transfers phase and history information through
mutable Fortran \texttt{COMMON} blocks and also contains \texttt{SAVE} and
\texttt{DATA} storage.  Threads in one process share these objects, so
concurrent UEL and UMAT calls can create read--write or write--write races
unless ownership is guaranteed or access is synchronized.  Separate MPI
processes present a different problem: each rank has its own common-block
instance, and common storage therefore cannot provide inter-rank exchange.
Distributed support would require an explicit ownership and communication
design at defined synchronization points.

The earlier D2C comparison produced identical serial and one-rank/four-thread
phase, history, reaction-force, energy, and increment results.  This is
positive but case-specific evidence; it does not prove general thread safety.
Stage~P therefore isolates the question in an eight-physical-element model.
Bounded diagnostics record callback identity, rank and thread identifiers,
first access to each shared location, initialization, live writer conflicts,
and final call counts.  A serial diagnostic is the mandatory reference before
one four-thread characterization.  MPI, hybrid, and production claims remain
withheld.

The first instrumented serial package terminated before element execution
because \texttt{GETRANK} was called during \texttt{UEXTERNALDB} initialization.
The subsequent uninstrumented P3-SB model compiled, linked, completed
Abaqus/Standard, and produced complete eight-element/four-integration-point
state evidence. Its automated lane nevertheless retained a technical
validation-failure classification because the required standard-output and
status files were copied only after validation. An isolated replay of the
unchanged validator against the finalized evidence passed every frozen serial
gate, including 32/32 state records, finite reaction-force and energy
histories, monotone phase/history fields, and 13 increment/attempt records.
This supports the eight-element serial baseline but no thread- or MPI-safety
claim. The subsequent P3-SM0 test used only constant \texttt{UEXTERNALDB},
UEL, and UMAT markers; rank/thread and mutex utilities remained excluded.
The one authorized serial job completed with scheduler and solver exit zero,
all four markers observed, no signal 11, complete 32/32 state coverage, and
zero phase, history, or transfer violations. Thus, the accepted serial model
completes with this minimal callback plumbing. The result does not establish
\texttt{GETRANK} or \texttt{GETTHREADID} safety, shared
\texttt{COMMON}/\texttt{SAVE} thread safety, four-thread repeatability, MPI,
or hybrid support.
\subsection{Separated identifier-utility qualification}
P3-SM0 passed serially with minimal \texttt{UEXTERNALDB}, UEL, and UMAT
callbacks, solver exit zero, complete 32-record state coverage, and no signal
11.  P3-SM1T is prepared, but not authorized for execution, to call only
\texttt{GETTHREADID()} inside the controlled UEL condition for element 1,
step 1, increment 1.  Markers immediately bracket every call, and repeated
pairs are counted.  P3-SM1R remains design-only; \texttt{GETRANK} remains
unqualified outside the failed \texttt{UEXTERNALDB(LOP=0)} location.  P3-T4
and all parallel or downstream production routes remain blocked.

The single P3-SM1T execution compiled and linked the user source and completed
input processing.  In the controlled UEL callback, the before marker was
written, but the Abaqus/Standard executable then failed dynamic symbol
resolution for \texttt{getthreadid\_}.  Consequently the after-marker count
was zero, no thread identifier was returned, and no signal 11 occurred.  This
is classified as an identifier-interface failure, not as evidence about
thread safety.

An offline audit then resolved the Abaqus 2023 installation headers.  The
installed \texttt{SMAAspUserSubroutines.hdr} includes a utility interface
declaring \texttt{get\_thread\_id()}, whereas it does not declare
\texttt{GETTHREADID()}.  Thus the closed job tested the undocumented spelling
that requested \texttt{getthreadid\_}, not the installed interface.  A new
P3-SM1TC package uses the installed spelling and remains unauthorized.  The
previous P3-S source likewise used GETRANK as a function expression rather
than the documented subroutine-call form; P3-SM1R remains design-only.

The separately authorized P3-SM1TC serial execution then passed.  Three
controlled UEL invocations returned thread identifier zero, with matched
before/after markers, no signal 11, and no unresolved symbol.  All P3-SM0
callback and scientific gates also passed.  The supported conclusion is
limited to the installed \texttt{get\_thread\_id()} interface under one MPI
rank and one OpenMP thread.  GETRANK and multi-thread behavior remain
unqualified.

The subsequent P3-SM1R execution also passed: three controlled UEL calls to
the documented \texttt{CALL GETRANK(KPROCESSNUM)} interface returned rank
zero, with complete serial scientific gates.  This result remains limited to
one MPI rank and one OpenMP thread.

P3-T4 is therefore prepared offline as a bounded one-rank, four-thread
characterization of the eight-element case.  It retains the P3-SM0 scientific
expressions and brackets their unsynchronized shared-state accesses with
diagnostic metadata calls.  Abaqus mutexes protect only counters, ownership,
and active-writer metadata; they do not protect the scientific
\texttt{USRVAR} or \texttt{TRANSFER\_DONE} operations.  Execution remains
unauthorized, and the preparation makes no MPI, hybrid, or general production
thread-safety claim.
